\documentclass[11pt]{article}

\usepackage[a4paper,margin=29mm]{geometry}
\usepackage[T1]{fontenc}
\usepackage[utf8]{inputenc}
\usepackage{lmodern}
\usepackage{microtype}
\usepackage{amsmath,amssymb,amsthm,mathtools}
\usepackage{booktabs,tabularx,array}
\usepackage{enumitem}
\usepackage{xcolor}
\usepackage[round,authoryear]{natbib}
\usepackage[colorlinks=true,linkcolor=blue!55!black,
  citecolor=red!45!black,urlcolor=blue!55!black]{hyperref}
\usepackage[nameinlink,noabbrev]{cleveref}

\hypersetup{
  pdftitle={Stable Marriage, Optimal Transport, and Birkhoff Geometry},
  pdfauthor={Gabriel Peyre},
  pdfsubject={Gale--Shapley stability, Kantorovich assignment, and the assignment game},
  pdfkeywords={stable marriage, Gale--Shapley, optimal transport, Birkhoff polytope, assignment game}
}

\newtheorem{theorem}{Theorem}[section]
\newtheorem{proposition}[theorem]{Proposition}
\newtheorem{lemma}[theorem]{Lemma}
\newtheorem{corollary}[theorem]{Corollary}
\theoremstyle{definition}
\newtheorem{definition}[theorem]{Definition}
\newtheorem{example}[theorem]{Example}
\theoremstyle{remark}
\newtheorem{remark}[theorem]{Remark}

\newcommand{\R}{\mathbb R}
\newcommand{\one}{\mathbf 1}
\newcommand{\Perm}{\mathfrak S}
\newcommand{\Pmat}{\mathrm P}
\newcommand{\Cmat}{\mathrm C}
\newcommand{\Smat}{\mathrm S}
\newcommand{\fD}{\mathrm f}
\newcommand{\gD}{\mathrm g}
\newcommand{\Bcal}{\mathcal B}
\newcommand{\Scal}{\mathcal S}
\newcommand{\Ucal}{\mathcal U}
\newcommand{\argmin}{\operatorname*{arg\,min}}
\newcommand{\argmax}{\operatorname*{arg\,max}}
\newcommand{\ip}[2]{\left\langle #1,#2\right\rangle}
\newcommand{\rank}{\operatorname{rank}}

\title{\textbf{Stable Marriage, Optimal Transport,\\
and Birkhoff Geometry}\\
\large Ordinal stability, cardinal efficiency, and transferable utility}
\author{Gabriel Peyr\'e}
\date{July 2026}

\begin{document}
\maketitle

\begin{abstract}
Stable marriage and discrete optimal transport both pair two finite populations,
but they answer different questions. Gale--Shapley stability is ordinal and
local: no unmatched pair should prefer one another to their assigned partners.
Kantorovich transport is cardinal and global: it minimizes an aggregate cost
over the Birkhoff polytope. This note separates these notions before drawing
their precise connections. The Birkhoff--von Neumann theorem identifies the
Birkhoff polytope with the convex hull of all perfect matchings, while the
Vande Vate--Rothblum theorem identifies a smaller polytope, obtained by adding
linear no-blocking inequalities, with the convex hull of stable matchings.
Ordinary stability is not equivalent to minimum total cost, even when both
sides rank partners using the same cost matrix. The exact optimal-transport
bridge instead appears in the Shapley--Shubik assignment game: with
transferable utility, pairwise stability is Kantorovich dual feasibility,
matched pairs satisfy complementary slackness, and Birkhoff integrality yields
an integral stable assignment. We also explain what Gale--Shapley optimizes,
compare its algorithmic invariants with Hungarian primal--dual methods, and
discuss unequal populations, fractional masses, and continuous matching.
\end{abstract}

\section{Three questions about the same pairing matrix}
\label{sec:three-questions}

Let $X=\{1,\ldots,n\}$ and $Y=\{1,\ldots,n\}$ be two populations. A
perfect matching is a permutation $\sigma\in\Perm_n$: agent $i\in X$ is
paired with \(\sigma(i)\in Y\). Its permutation matrix is
\begin{equation}\label{eq:permutation-matrix}
  (\Pmat_\sigma)_{ij}=\begin{cases}1,&j=\sigma(i),\\0,&\text{otherwise}.
  \end{cases}
\end{equation}
This common representation hides three different mathematical models.

\begin{enumerate}[label=\textbf{\arabic*.},leftmargin=2.7em]
  \item \textbf{Optimal assignment.} A cardinal cost $C_{ij}$ measures the
  loss incurred by matching $i$ to $j$. One minimizes the sum of costs.
  \item \textbf{Stable marriage.} Each agent only ranks the agents on the
  opposite side. One asks that no unmatched pair mutually prefer one another
  to their assigned partners.
  \item \textbf{Assignment game.} A match produces a cardinal surplus that
  the two partners can divide through a transfer. One asks that no pair can
  deviate and split its surplus so as to improve both partners' payoffs.
\end{enumerate}

The first model is a social planner's optimization problem, the second is a
no-blocking equilibrium without transfers, and the third is a market
equilibrium with transferable utility. Confusing them leads to false claims.
In particular, a stable marriage need not minimize any natural aggregate cost
on all matchings. The assignment game, not ordinary stable marriage, is the
setting in which Kantorovich duality becomes an exact stability theorem.

\section{Optimal assignment and the Birkhoff polytope}
\label{sec:birkhoff}

We first recall the convex geometry underlying discrete transport. The
Birkhoff polytope is
\begin{equation}\label{eq:birkhoff-polytope}
  \Bcal_n
  =\left\{\Pmat\in\R_+^{n\times n}:
  \Pmat\one=\one,\ \Pmat^\top\one=\one\right\}.
\end{equation}
Its elements are bistochastic matrices. If both empirical measures assign mass
\(1/n\) to every agent, the corresponding Kantorovich couplings are
$\Pmat/n$ with $\Pmat\in\Bcal_n$. We retain the unnormalized convention
\eqref{eq:birkhoff-polytope} because its integral points are exactly the
permutation matrices.

\begin{definition}[Optimal assignment and its Kantorovich relaxation]
For a cost matrix $\Cmat=(C_{ij})\in\R^{n\times n}$, define
\begin{align}
  \operatorname{Assign}(\Cmat)
  &=\min_{\sigma\in\Perm_n}\sum_{i=1}^n C_{i,\sigma(i)},
  \label{eq:assignment}\\
  \operatorname{Kant}(\Cmat)
  &=\min_{\Pmat\in\Bcal_n}\ip{\Cmat}{\Pmat}
  =\min_{\Pmat\in\Bcal_n}\sum_{i,j}C_{ij}P_{ij}.
  \label{eq:kantorovich-assignment}
\end{align}
\end{definition}

\begin{theorem}[Birkhoff--von Neumann]\label{thm:birkhoff}
The extreme points of $\Bcal_n$ are the permutation matrices. Equivalently,
every $\Pmat\in\Bcal_n$ admits a decomposition
\begin{equation}\label{eq:birkhoff-decomposition}
  \Pmat=\sum_{k=1}^N\lambda_k\Pmat_{\sigma_k},
  \qquad \lambda_k\geq0,\quad \sum_{k=1}^N\lambda_k=1.
\end{equation}
Consequently $\operatorname{Kant}(\Cmat)=\operatorname{Assign}(\Cmat)$,
and the Kantorovich problem has an integral minimizer.
\end{theorem}

\begin{proof}
The support graph of a bistochastic matrix has no degree-one vertex unless its
incident entry is one. If a nonintegral bistochastic matrix has no such forced
edge, its bipartite support contains an even cycle. Alternately adding and
subtracting a sufficiently small amount along that cycle writes the matrix as
the midpoint of two distinct bistochastic matrices, so it is not extreme.
Conversely, a permutation matrix cannot be written as a nontrivial convex
combination of bistochastic matrices because each of its zero entries forces
the corresponding entries of both matrices to vanish. This proves the
extreme-point statement and hence \eqref{eq:birkhoff-decomposition}. A linear
functional reaches its minimum at an extreme point, proving the last claim.
See \citet{Birkhoff1946,VonNeumann1953}.
\end{proof}

The linear-programming dual of \eqref{eq:kantorovich-assignment} is
\begin{equation}\label{eq:kantorovich-dual}
  \max_{\fD,\gD\in\R^n}
  \left\{\sum_i f_i+\sum_j g_j:
  f_i+g_j\leq C_{ij}\ \text{for all }i,j\right\}.
\end{equation}
At optimality, complementary slackness gives
\begin{equation}\label{eq:complementary-slackness-cost}
  P_{ij}>0\quad\Longrightarrow\quad f_i+g_j=C_{ij}.
\end{equation}
The Hungarian algorithm is a primal--dual method built around precisely these
feasible potentials and equality edges \citep{Kuhn1955}.

\section{Stable marriage and deferred acceptance}
\label{sec:stable-marriage}

Stable marriage replaces cardinal costs by ordinal preferences. For every
\(i\in X\), let \(\succ_i^X\) be a strict total order on \(Y\);
\(j\succ_i^X j'\) means that \(i\) prefers \(j\) to \(j'\). For every
\(j\in Y\), let \(\succ_j^Y\) be a strict total order on \(X\).

\begin{definition}[Blocking pair and stable matching]\label{def:stable}
For a matching \(\sigma\in\Perm_n\), an unmatched pair \((i,j)\), with
\(j\neq\sigma(i)\), is a \emph{blocking pair} when
\begin{equation}\label{eq:blocking-pair}
  j\succ_i^X\sigma(i)
  \qquad\text{and}\qquad
  i\succ_j^Y\sigma^{-1}(j).
\end{equation}
The matching is \emph{stable} if it has no blocking pair.
\end{definition}

Stability is local and ordinal. It does not sum gains across agents, and a
large improvement for one pair cannot compensate for a loss to another pair.
Gale and Shapley's deferred-acceptance algorithm proves that a stable matching
always exists under strict complete preferences \citep{GaleShapley1962}.

\paragraph{The \(X\)-proposing deferred-acceptance algorithm.}
Initially all agents are unmatched. While an unmatched \(i\in X\) has not
proposed to everyone, \(i\) proposes to the most-preferred \(j\in Y\) who has
not yet rejected him. Each \(j\) tentatively holds her most-preferred proposal
among her current tentative partner and all new proposals, and rejects the
others. The tentative matches become final when no further proposal is
possible.

\begin{theorem}[Gale--Shapley]\label{thm:gale-shapley}
With strict complete preferences, deferred acceptance terminates after at most
\(n^2\) proposals and returns a perfect stable matching \(\sigma_X\). Moreover,
it is \(X\)-optimal among stable matchings: for every stable \(\sigma\) and
every \(i\in X\),
\begin{equation}\label{eq:proposer-optimal}
  \sigma_X(i)\succeq_i^X\sigma(i).
\end{equation}
If the roles are reversed, the resulting matching is \(Y\)-optimal.
\end{theorem}

\begin{proof}
No ordered pair \((i,j)\) receives more than one proposal, so there are at most
\(n^2\) proposals. Once a receiver holds a proposal, she always holds one and
her tentative partner can only improve. At termination every proposer is
matched. Indeed, if some proposer had exhausted all receivers, every receiver
would have received and hence would still hold one proposal, accounting for
all \(n\) proposers.

To prove stability, suppose that \((i,j)\) blocks the output. Since \(i\)
prefers \(j\) to his final partner, he proposed to \(j\) earlier. Receiver
\(j\) rejected him while holding someone she preferred, and her tentative
partner only improved afterward. She therefore prefers her final partner to
\(i\), contradicting \eqref{eq:blocking-pair}.

For proposer optimality, suppose for contradiction that some proposer is
rejected by a receiver with whom he is paired in some stable matching. Consider
the first such rejection: \(j\) rejects \(i\) while holding \(i'\), whom she
prefers to \(i\). In a stable matching \(\sigma\) with \(\sigma(i)=j\), write
\(j'=\sigma(i')\). Proposer \(i'\) must prefer \(j\) to \(j'\). Otherwise he
would already have proposed to, and been rejected by, his stable partner
\(j'\), contradicting the choice of the first rejection. Thus \(i'\) prefers
\(j\) to \(j'\), while \(j\) prefers \(i'\) to \(i\), so \((i',j)\) blocks
\(\sigma\). This contradiction shows that no proposer is ever rejected by a
possible stable partner. His final partner is therefore his best stable
partner, proving \eqref{eq:proposer-optimal}. Reversing the sides gives the
last claim.
\end{proof}

Stable matchings form a distributive lattice under the order induced by the
\(X\)-agents' preferences; the two deferred-acceptance outcomes are its
opposite extrema. This finer structure is developed by
\citet{GusfieldIrving1989} and \citet{RothSotomayor1990}.

\section{A stable analogue of the Birkhoff polytope}
\label{sec:stable-polytope}

The parallel with Kantorovich transport becomes geometric once stability is
written as linear inequalities. For a pair \((i,j)\), set
\begin{equation}\label{eq:better-sets}
  B_i(j)=\{j'\in Y:j'\succ_i^X j\},
  \qquad
  B_j(i)=\{i'\in X:i'\succ_j^Y i\}.
\end{equation}
Define the fractional stable-marriage polytope
\begin{equation}\label{eq:stable-polytope}
  \Scal(\succ^X,\succ^Y)
  =\left\{\Pmat\in\Bcal_n:
  P_{ij}+\sum_{j'\in B_i(j)}P_{ij'}
  +\sum_{i'\in B_j(i)}P_{i'j}\geq1
  \quad\forall(i,j)\right\}.
\end{equation}

\begin{proposition}[Integral no-blocking inequalities]\label{prop:integral-stability}
For \(\sigma\in\Perm_n\), its permutation matrix \(\Pmat_\sigma\) belongs to
\(\Scal(\succ^X,\succ^Y)\) if and only if \(\sigma\) is stable.
\end{proposition}

\begin{proof}
For an integral matching, the left-hand side of the inequality indexed by
\((i,j)\) is positive precisely when at least one of the following holds:
\(i\) is matched to \(j\); \(i\) is matched to someone preferred to \(j\); or
\(j\) is matched to someone preferred to \(i\). The inequality therefore fails
exactly when \(i\) and \(j\) are unmatched and both prefer one another to their
current partners, which is exactly the blocking condition.
\end{proof}

The difficult statement is that the inequalities introduce no fractional
extreme points. This is the stable-marriage counterpart of the
Birkhoff--von Neumann theorem.

\begin{theorem}[Vande Vate--Rothblum stable-polytope theorem]
\label{thm:stable-polytope}
Under strict complete preferences,
\begin{equation}\label{eq:stable-convex-hull}
  \Scal(\succ^X,\succ^Y)
  =\operatorname{conv}\left\{\Pmat_\sigma:
  \sigma\text{ is stable}\right\}.
\end{equation}
In particular, every extreme point of \eqref{eq:stable-polytope} is the
permutation matrix of a stable matching.
\end{theorem}

\begin{proof}[Proof structure]
Proposition~\ref{prop:integral-stability} gives the inclusion from right to
left. The reverse inclusion is the nontrivial polyhedral theorem of
\citet{VandeVate1989}, simplified and extended by \citet{Rothblum1992}.
Rothblum proves that every linear objective over the system
\eqref{eq:stable-polytope} has an integral optimum. One can formulate the key
step through the rotation poset of stable matchings: eliminating a closed set
of rotations moves from the \(X\)-optimal matching to another stable matching,
and the resulting order-polytope representation is integral. Equivalently, a
linear image of that integral order polytope gives
\eqref{eq:stable-polytope}. Later, \citet{KiralPap2008} proved that Rothblum's
inequality system is totally dual integral.

For completeness, once this integrality result is granted, the convex-hull
claim follows immediately. If \eqref{eq:stable-polytope} had a fractional
extreme point \(\Pmat\), choose a linear functional that uniquely exposes it.
The integrality theorem would supply an integral optimizer for the same
functional, contradicting uniqueness. Hence every extreme point is integral,
and Proposition~\ref{prop:integral-stability} identifies it with a stable
matching. A bounded polytope is the convex hull of its extreme points.
\end{proof}

\begin{remark}[Two different integrality mechanisms]
Theorem~\ref{thm:birkhoff} says that every bistochastic matrix is a mixture of
arbitrary perfect matchings. Theorem~\ref{thm:stable-polytope} says that every
bistochastic matrix satisfying the fractional no-blocking inequalities is a
mixture of stable matchings. An arbitrary Birkhoff decomposition of a point in
\(\Scal\) need not consist of stable permutations; the theorem guarantees that
at least one stable decomposition exists.
\end{remark}

\begin{corollary}[Minimum-cost stable marriage]\label{cor:min-cost-stable}
For every weight matrix \(\Cmat\), the linear program
\begin{equation}\label{eq:min-cost-stable}
  \min_{\Pmat\in\Scal(\succ^X,\succ^Y)}\ip{\Cmat}{\Pmat}
\end{equation}
has an integral optimum and therefore computes a minimum-cost stable matching.
\end{corollary}

This looks like Kantorovich transport, but the feasible polytope now depends on
the two preference profiles. The cost only selects among stable matchings; it
does not create stability.

\paragraph{What deferred acceptance optimizes.}
Let \(r^X_{ij}\in\{1,\ldots,n\}\) be the rank of \(j\) in proposer \(i\)'s
preference list, with rank one best. The pointwise optimality
\eqref{eq:proposer-optimal} is stronger than minimizing one aggregate score.

\begin{corollary}[A linear-programming characterization of Gale--Shapley]
\label{cor:gs-lp}
For any weights \(w_i>0\), the \(X\)-proposing Gale--Shapley output solves
\begin{equation}\label{eq:gs-rank-objective}
  \min_{\Pmat\in\Scal(\succ^X,\succ^Y)}
  \sum_{i,j}w_i r^X_{ij}P_{ij}.
\end{equation}
It simultaneously minimizes every proposer's rank among stable matchings.
\end{corollary}

\begin{proof}
By Theorem~\ref{thm:stable-polytope}, every feasible \(\Pmat\) is a convex
combination of stable permutation matrices. For each such matching,
\eqref{eq:proposer-optimal} gives
\(r^X_{i,\sigma_X(i)}\leq r^X_{i,\sigma(i)}\) for every \(i\). Multiplication by
\(w_i>0\), summation, and convex combination prove the claim.
\end{proof}

This does not imply that the \(X\)-optimal matching minimizes a two-sided rank
sum or a cardinal transport cost over all matchings.

\section{Why ordinary stability is not optimal transport}
\label{sec:not-equivalent}

It is tempting to generate both preference profiles from a common cost matrix:
each agent prefers the partner with smaller cost. Even then, no-blocking
stability and minimum total cost generally disagree.

\begin{example}[A stable matching can be globally expensive]
\label{ex:counterexample}
Consider
\begin{equation}\label{eq:counterexample-cost}
  \Cmat=\begin{pmatrix}0&1\\50&100\end{pmatrix},
\end{equation}
and let every row and every column rank partners by increasing \(C_{ij}\). The
diagonal matching has total cost \(100\) and is stable. The pair \((1,2)\) does
not block because agent \(1\in X\) prefers \(1\in Y\), while \((2,1)\) does not
block because agent \(1\in Y\) prefers \(1\in X\). The off-diagonal matching
has the smaller total cost \(1+50=51\), so it is the unique optimal assignment,
but the pair \((1,1)\) blocks it. Thus the unique stable matching differs from
the unique Kantorovich optimum.
\end{example}

There is also a geometric obstruction to encoding all stable matchings as the
minimizers of a single unconstrained transport cost. The minimizers of a linear
functional over \(\Bcal_n\) form a face of \(\Bcal_n\). The convex hull of the
stable matchings is generally a polytope cut out by additional inequalities,
not a face of \(\Bcal_n\). Hence there need not exist a matrix \(\Cmat\) whose
entire set of optimal assignments is exactly the set of stable marriages.

The two notions answer different welfare questions:
\begin{itemize}[leftmargin=2em]
  \item optimal transport permits losses and gains to compensate across
  unrelated agents through the sum \(\sum_i C_{i,\sigma(i)}\);
  \item stable marriage forbids a deviation only when both members of one pair
  improve, and it allows no monetary compensation between them.
\end{itemize}

\section{The exact bridge: transferable utility and the assignment game}
\label{sec:assignment-game}

The connection with Kantorovich duality becomes exact once utility is
transferable. Let \(S_{ij}\) be the total surplus created by matching \(i\) and
\(j\). For simplicity assume that all agents are matched and that outside
options have been normalized to zero. A matched pair may divide \(S_{ij}\)
arbitrarily: \(u_i\) is the payoff of \(i\in X\), and \(v_j\) that of \(j\in Y\).

\begin{definition}[Core-stable assignment-game outcome]
An outcome \((\sigma,u,v)\) is core stable when
\begin{align}
  u_i+v_{\sigma(i)}&=S_{i,\sigma(i)} &&\text{for every }i,
  \label{eq:budget-balance}\\
  u_i+v_j&\geq S_{ij} &&\text{for every }(i,j),
  \label{eq:no-transfer-blocking}
\end{align}
with \(u_i,v_j\geq0\) when individual rationality relative to the outside
option is required.
\end{definition}

Indeed, pair \((i,j)\) can make both members strictly better through a transfer
if and only if its available surplus satisfies \(S_{ij}>u_i+v_j\). Thus
\eqref{eq:no-transfer-blocking} is exactly the no-blocking condition with
transfers.

Consider the primal and dual linear programs
\begin{align}
  \max_{\Pmat\in\Bcal_n}\ip{\Smat}{\Pmat},
  \label{eq:surplus-primal}\\
  \min_{u,v\in\R^n}
  \left\{\sum_i u_i+\sum_jv_j:
  u_i+v_j\geq S_{ij}\quad\forall(i,j)\right\}.
  \label{eq:surplus-dual}
\end{align}

\begin{theorem}[Shapley--Shubik: core stability is Kantorovich duality]
\label{thm:assignment-game}
An outcome \((\sigma,u,v)\) is core stable if and only if
\(\Pmat_\sigma\) solves \eqref{eq:surplus-primal}, \((u,v)\) solves
\eqref{eq:surplus-dual}, and the two solutions satisfy complementary slackness.
Consequently the assignment game has a core-stable outcome.
\end{theorem}

\begin{proof}
Suppose first that the outcome is core stable. The inequalities
\eqref{eq:no-transfer-blocking} make \((u,v)\) dual feasible. Summing
\eqref{eq:budget-balance} over matched pairs yields
\begin{equation}\label{eq:core-equal-values}
  \ip{\Smat}{\Pmat_\sigma}
  =\sum_i S_{i,\sigma(i)}
  =\sum_i u_i+\sum_jv_j.
\end{equation}
Weak duality says that every primal value is at most every dual value, so
equality proves that both candidates are optimal. Equality on matched pairs is
complementary slackness.

Conversely, let \(\Pmat_\sigma\) and \((u,v)\) be optimal primal and dual
solutions. Strong linear-programming duality equates their objectives.
Complementary slackness gives
\(u_i+v_{\sigma(i)}=S_{i,\sigma(i)}\), while dual feasibility gives
\(u_i+v_j\geq S_{ij}\) for every pair. These are precisely
\eqref{eq:budget-balance}--\eqref{eq:no-transfer-blocking}.

Finally, linear-programming duality supplies a dual optimizer, and
Theorem~\ref{thm:birkhoff} supplies an integral primal optimizer. Together they
form a core-stable outcome. This is the finite assignment-game theorem of
\citet{ShapleyShubik1971}.
\end{proof}

To recover the cost convention, set \(S_{ij}=-C_{ij}\), \(f_i=-u_i\), and
\(g_j=-v_j\). Then \eqref{eq:surplus-primal} becomes the Kantorovich
minimization \eqref{eq:kantorovich-assignment}, while
\(u_i+v_j\geq S_{ij}\) becomes the familiar dual inequality
\(f_i+g_j\leq C_{ij}\) from \eqref{eq:kantorovich-dual}. Equality on matched
pairs becomes \eqref{eq:complementary-slackness-cost}.

\begin{remark}[Why transfers change the mathematics]
In ordinary stable marriage, the two comparisons in
\eqref{eq:blocking-pair} are separate ordinal statements. In the assignment
game, a transfer lets the pair pool its surplus, replacing two preference
comparisons by the scalar inequality \(u_i+v_j\geq S_{ij}\). That inequality
is a Kantorovich dual constraint. This aggregation is the exact reason optimal
transport enters the transferable-utility model but not the general
Gale--Shapley model.
\end{remark}

The continuous counterpart is equally important in economics. Quasilinear
hedonic equilibria and stable matching with transferable utility can be cast as
Monge--Kantorovich problems; see \citet{ChiapporiMcCannNesheim2010} for
existence, purity, and uniqueness questions in continuous type spaces.

\section{Three integrality and duality mechanisms}
\label{sec:dictionary}

The comparison can be summarized without collapsing the models into one
another.

\begin{center}
\small
\renewcommand{\arraystretch}{1.35}
\begin{tabularx}{\textwidth}{>{\raggedright\arraybackslash}p{0.18\textwidth}
  >{\raggedright\arraybackslash}p{0.24\textwidth}
  >{\raggedright\arraybackslash}p{0.26\textwidth}
  >{\raggedright\arraybackslash}X}
\toprule
Model & Data and criterion & Convex geometry & Main algorithmic principle \\
\midrule
Optimal assignment & Cardinal costs \(C_{ij}\); minimize total cost &
\(\Bcal_n=\operatorname{conv}\{\Pmat_\sigma\}\) by Birkhoff--von Neumann &
Hungarian or auction primal--dual optimization \\
Stable marriage & Two ordinal preference profiles; exclude blocking pairs &
\(\Scal=\operatorname{conv}\{\Pmat_\sigma:\sigma\text{ stable}\}\) by
Vande Vate--Rothblum & Deferred acceptance; irreversible rejections and
one-sided lattice optimum \\
Assignment game & Cardinal surplus \(S_{ij}\) with transfers; exclude
profitable pair deviations & Primal Birkhoff polytope plus dual payoff
inequalities \(u_i+v_j\geq S_{ij}\) & Linear-programming duality,
complementary slackness, and market-clearing prices \\
\bottomrule
\end{tabularx}
\end{center}

The words \emph{primal} and \emph{dual} play different roles here. In the
stable-marriage polytope, no-blocking is imposed by additional primal
inequalities on \(\Pmat\). In the assignment game, no-blocking is exactly dual
feasibility for the payoffs. The latter is why Hungarian potentials can be
read as equilibrium utilities or prices. Gale--Shapley's tentative partners
and rejections are not Kantorovich potentials, even though both algorithms
maintain certificates that monotonically restrict the remaining choices.

\section{Extensions and limits of the analogy}
\label{sec:extensions}

\paragraph{Unequal populations and outside options.}
For optimal transport, unequal cardinalities pose no conceptual difficulty:
one uses rectangular coupling matrices with prescribed marginals. For stable
marriage, one allows unmatched agents and declares being unmatched an outside
option, or adds dummy agents. In the assignment game, dummy partners with zero
surplus produce the same construction and make individual rationality explicit.

\paragraph{Many-to-one matching and capacities.}
College admissions replaces one-to-one matching by capacities on the receiving
side. Deferred acceptance still applies under responsive preferences. Its
feasible incidence matrices resemble a capacitated transportation polytope,
but stability again adds preference-dependent no-blocking constraints. The
transportation constraints and the stability constraints should not be
identified.

\paragraph{Fractional masses.}
Kantorovich transport naturally splits mass and accommodates arbitrary
histograms. Stable marriage is instead an indivisible-agent model. Theorem
\ref{thm:stable-polytope} permits a fractional point only as a lottery over
integral stable marriages; its integrality is special to bipartite stable
matching with the stated preference assumptions. With ties, couples,
nonbipartite roommates, or more general complementarities, existence and
integrality can fail or require different stability notions.

\paragraph{Randomization.}
A bistochastic matrix can be interpreted as marginals of a lottery over
matchings. Birkhoff--von Neumann says every such fractional assignment can be
implemented by randomizing over deterministic assignments. If the matrix also
lies in \(\Scal\), Vande Vate--Rothblum says one can randomize using only stable
assignments. This is stronger than checking expected preference scores: ex
post every realized matching can be stable.

\paragraph{Continuous matching.}
In continuous populations, optimal transport retains its measure-theoretic
meaning. Stable matching without transfers needs a suitable continuum market
model and is not simply a transport problem. Under transferable utility and
quasilinear preferences, however, Kantorovich dual potentials become indirect
utilities or prices, and complementary slackness identifies equilibrium
matches. This is the robust path from finite assignment games to continuous
optimal transport.

\section{Takeaway}

The clean connection consists of two parallels and one exact equivalence.
\begin{enumerate}[label=\arabic*.,leftmargin=2em]
  \item \textbf{Parallel convex hulls.} Birkhoff--von Neumann convexifies all
  perfect matchings; Vande Vate--Rothblum convexifies the stable ones after
  adding linear no-blocking inequalities.
  \item \textbf{Different objectives.} Ordinary stable marriage is an ordinal
  equilibrium problem, whereas Kantorovich assignment is a cardinal welfare
  problem. Neither generally implies the other.
  \item \textbf{Exact transferable-utility bridge.} In the assignment game,
  core stability is Kantorovich dual feasibility, matched pairs are contact
  points of the dual constraint, and complementary slackness plus Birkhoff
  integrality produces an integral stable allocation.
\end{enumerate}
Thus Gale--Shapley and optimal transport belong to the same broad geometry of
bipartite couplings, but the assignment game is the precise point where their
notions of stability and optimality coincide.

\bibliographystyle{plainnat}
\bibliography{references}

\end{document}
